\section{Frenet--Serret Frame and Formulas}

\begin{enumerate}
  \item What is a regular curve? What condition must a parametrized curve $\mathbf{r}(t)$ satisfy?
  \item What is arc-length parametrization? How do we define the arc-length parameter $s$?
  \item Why is arc-length parametrization especially convenient when defining the Frenet frame?
  \item How is the unit tangent vector $\mathbf{T}(s)$ defined in terms of $\mathbf{r}(s)$?
  \item Why does $\|\mathbf{T}(s)\| = 1$ always hold?
  \item Show that $\mathbf{T}'(s)$ is orthogonal to $\mathbf{T}(s)$. Why must this be true?
  \item How is curvature $\kappa(s)$ defined in terms of $\mathbf{T}'(s)$?
  \item What is the geometric meaning of curvature? What does it measure intuitively?
  \item Under what condition is curvature zero? What does this imply about the curve?
  \item How is the principal normal vector $\mathbf{N}(s)$ defined?
  \item Why is $\mathbf{N}(s)$ a unit vector?
  \item What is the binormal vector $\mathbf{B}(s)$ and how is it defined?
  \item Why is $\mathbf{B}(s)$ orthogonal to both $\mathbf{T}(s)$ and $\mathbf{N}(s)$?
  \item Show that $\{\mathbf{T}, \mathbf{N}, \mathbf{B}\}$ forms an orthonormal basis.
  \item What is the geometric interpretation of the Frenet frame at a point on the curve?
  \item Write down the Frenet--Serret formulas.
  \item What does the equation $\mathbf{T}'(s) = \kappa(s)\mathbf{N}(s)$ tell us geometrically?
  \item How is torsion $\tau(s)$ defined?
  \item What is the geometric meaning of torsion?
  \item What does it mean if torsion is zero everywhere?
  \item Explain the equation $\mathbf{B}'(s) = -\tau(s)\mathbf{N}(s)$.
  \item Why must $\mathbf{N}'(s)$ involve both $\mathbf{T}(s)$ and $\mathbf{B}(s)$?
  \item Derive the Frenet--Serret formulas starting from orthonormality constraints.
  \item How do curvature and torsion uniquely determine a curve (up to rigid motion)?
  \item What is the fundamental theorem of space curves?
  \item Compute $\kappa$ and $\tau$ for a straight line.
  \item Compute $\kappa$ and $\tau$ for a circle of radius $R$.
  \item Compute $\kappa$ and $\tau$ for a helix.
  \item How do you compute curvature for a curve given in a non arc-length parameter $t$?
  \item What is the formula for curvature using $\mathbf{r}'(t)$ and $\mathbf{r}''(t)$?
  \item How is torsion computed using $\mathbf{r}'(t)$, $\mathbf{r}''(t)$, and $\mathbf{r}'''(t)$?
  \item Why does the cross product $\mathbf{r}'(t) \times \mathbf{r}''(t)$ appear in curvature and torsion formulas?
  \item What happens to the Frenet frame at points where $\kappa = 0$?
  \item Can the Frenet frame always be defined globally along a curve?
  \item What is the osculating plane? Which vectors span it?
  \item What is the normal plane? Which vectors span it?
  \item What is the rectifying plane? Which vectors span it?
  \item How does the Frenet frame describe the local geometry of a curve?
  \item In what sense does curvature measure bending and torsion measure twisting?
  \item How would you reconstruct a curve from given curvature and torsion functions?
  \item What role do initial conditions play in reconstructing a curve?
  \item How does the Frenet frame evolve along the curve?
  \item What is the matrix form of the Frenet--Serret equations?
  \item Why is the Frenet--Serret matrix skew-symmetric (up to signs)?
  \item How is the Frenet frame related to rigid body motion along a curve?
  \item What is the connection between angular velocity and the Frenet--Serret formulas?
  \item How can the Frenet frame be interpreted as a moving coordinate system?
  \item What breaks down if the curve is not twice differentiable?
\end{enumerate}
